\chapter{Elementos do Modelo}
Como este modelo tem a pretenção de ser um metamodelo, este capítulo traz informações sobre sua estrutura bem como detalhes de seu uso. Nas próximas seções serão discutidas: a estrutura, a forma de imprimir a monografia, a ficha catalográfica, os índices, sumário, lista de acrônimos e bibliografia.

\section{Estrutura}
O modelo está dividido em três partes, que são: elementos pré-textuais, elementos textuais e 
elementos pós-textuais. A tabela \ref{estrutura} mostra essa estrutura.


\begin{table}  % aqui começa o ambiente tabela
\setlength{\baselineskip}{0.8\baselineskip} 
\centering
\caption{Estrutura do modelo}
\begin{tabular}{|c|l|}
\hline 
\hline
\textbf{Tipo }& \textbf{Elementos} \\ 
\hline 
\hline
\emph{Pré-textuais} & Capa\\ 
& Folha de rosto\\
& Ficha catalografica\\
& Folha de aprovacao\\
& Dedicatoria\\
& Agradecimentos\\
& Epígrafe\\
& Resumo\\
& Lista de figuras\\
& Lista de tabelas\\
& Sumário\\
& Lista de acrônimos 		\\
\hline 
\emph{Textuais} & Introdução \\ 
& Capítulo 2\\
& Capítulo 3\\
& Capítulo 4\\
& ...\\
& Conclusão \\
\hline 
\emph{Pós-textuais} & Referências\\
& Apêndices\\
& Anexos\\
\hline 

\end{tabular} \label{estrutura}

\end{table}

Todos\footnote{Exceto os índices, o sumário e a lista de acrônimos que são gerados automaticamente pelo \LaTeX} estes elementos estão separados em arquivos e estes em pastas, organizados conforme a tabela \ref{estrutura}:

\begin{description}
\item[pasta pretextuais] - capa.tex, 
folhaderosto.tex,
fichacatalografica.tex,
folhadeaprovacao.tex,
dedicatoria.tex,
agradecimentos.tex,
epigrafe.tex,
resumo.tex,
lista de figuras.tex,
lista de tabelas.tex,
sumário.tex,
lista de acronimos.tex
\item[pasta textuais] - introducao.tex, cap2.tex, cap3.tex, cap4.tex e conclusao.tex.
\item[pasta postextuais] - bibliografia.bib, apendice.tex e anexo.tex.
\end{description}


\nota{Para confeccionar o trabalho é necessário entrar em cada um destes arquivos e editá-los, com excessão dos arquivos gerados automaticamente que são: capa, folha de rosto e folha de aprovação. Sumário e listas também são geradas automaticamente.}


%\section{Impressão da Monografia}
%
%Este modelo está pré-configurado para imprimir em apenas um lado da folha, mas com uma alteração bem pequena é possível configurá-lo para impressão frente e verso. No preâmbulo do arquivo, a linha que define o tipo de documento é: 
%\begin{verbatim}
%\documentclass[12pt,a4paper]{abnt}
%\end{verbatim} 
%
%Para imprimir frente e verso substitua por:
%
%\begin{verbatim}
%\documentclass[12pt,a4paper,twoside]{abnt}
%\end{verbatim} 
%
%\nota{note que proposiltamente algumas folhas são deixadas em branco na impressão frente e verso: verso de sumário, verso de lista de figuras e verso de fim de capítulo.}
%
%É necessário descomentar a linha da contracapa para que o verso da capa fique em branco. E não se esqueça de configurar a impressora para impressão nos dois lados da folha.

\section{Ficha catalográfica}

Verifique junto a coordenação de curso a necessidade da ficha catalográfica. Caso seja necessária, a ficha catalográfica só deverá ser solicitada após a defesa e mediante a aprovação pela banca.  

Adquirida a ficha, salve-a no formato pdf com o nome de ficha.pdf, dentro da pasta figuras, vá ao arquivo fichacatalografica.tex (dentro da pasta pretextuais) e remova o símbolo de \% das linhas especificadas.

\section{Índices e sumário}
 	 	
Este modelo conta com um sumário, nomeado de ``conteúdo'' - conforme ABN\TeX, e também com índice de figuras e índice de tabelas. Destes, somente o sumário é obrigatório, os outros índices devem ser usados conforme a necessidade do trabalho, por exemplo, não é necessário inserir índice de tabelas se a monografia contém apenas duas tabelas. Lembrando que outros índices podem ser adicionados (de algoritmos).

O sumário e os índices de figuras e de tabelas são gerados automaticamente não é necessário editar, apenas compile-atualize o documento, que o \LaTeX\ faz o resto.

\nota{É importante lembrar que  mudanças aparecem depois de Compilar-atualizar duas ou três vezes.} 
 	
Para atualizar \verb+[no texmaker 3.5.2]+ ir em 

\verb+<Editar>-opção <Atualizar Estrutura>+


\section{Lista de acrônimos}

Para adicionar siglas use o comando \verb+ \sigla{SIGLA}{DESCRIÇÃO}+. Isso pode ser feito em qualquer parte do texto! Este modelo já está preparado para inserir a lista de acrônimos, contudo, é necessário gerá-la, para isso, vá ao terminal (no Texmaker o terminal está em \verb+ <Ferramentas>-opção <Abrir Terminal>+)	e digite a linha:


\begin{verbatim} makeindex monografia.nlo -s nomencl.ist -o monografia.nls 
\end{verbatim}


 \nota{Caso possua siglas e acrônimos, lembre se sempre de gerar a lista através do terminal}


\section{Bibliografia}

Edite o arquivo bibliografia, apenas acrescentando novas obras, não é necessário remover nada pois o latex insere somente as obras referenciadas. Ele também ordena tudo!!! A configuração padrão deste modelo é a ordem alfabética, uma alternativa seria a ordem de citação no texto.

Para inserir obras, utilize o BibTex ou BibLaTeX, preencha o maior número de campos possível.
O nome do autor é sempre importante, mesmo que a obra seja uma página da internet. O arquivo bibliografia.bib na pasta postextuais já conta com várias obras, de tipos bastante variados, não faz diferença removê-los ou não, a menos que sejam referenciadas em seu trabalho. 
A única recomendação é que utilize um arquivo chamado bibliografia.bib na pasta postextuais.
 
\subsection{Gerando as referências bibliográficas e citações}


Para gerar a bibliografia, devem ser realizados dois passos:

1. Compilar monografia.tex com o BibTex ( para geração do arquivo .bbl).

2. Atualizar a bibliografia (isso irá inserir as citações no texto).

Procure o manual do seu editor para saber como proceder. Exemplo: \verb+TeXMaker 3.3.3+
utilize:

1. F11 para gerar o arquivo .bbl e,

2. \verb+ menu <Editar>-opção <Atualizar Bibliografia>+ para atualizar a bibliografia no documento.

\nota{ É necessário compilar e atualizar o documento todas as vezes que inserir um novo item na bibliografia. Quase sempre será preciso compilar e atualizar duas ou três vezes para que as modificações sejam percebidas.} 



\section{Considerações finais do capítulo}

Este capítulo discutiu características de ....


No próximo capítulo serão abordados trabalhos que condizem com estado da arte em ....